\section{ Extensions and Computational Complexity}

\subsection{Model Extensions}

Our bicriteria results extend to the settings of repeated games, weighted voters, and multiple delegates.

\noindent $\bullet$ {\em Repeated Games}.
Recollect that an underlying motivation for liquid democracy is that it can be applied repeatedly over time with agents
having the option of delegating to different agents at different times. Evidently, by repeating the delegation
game over time, with different utility functions for the topics considered in different time periods, 
the same bicriteria guarantees holds.

\noindent $\bullet$ {\em Weighted Voters}. In some electoral systems, different agents may have different voting powers.
That is, the voters are weighted. In this case, the number of votes a guru casts is simply the sum of the weights of all the
votes to which it was delegated. Our bicriteria guarantees then follow trivially.

\noindent $\bullet$ {\em Multiple Delegates}.
In some settings it may be the case that an agent is allowed to nominate more than one delegate.
That are two natural models for this. First, an agent nominates multiple delegates but the mechanism
can use only one of them. Second, an agent splits its voting weight up and assigns it to multiple delegates.
In both cases our techniques can be adapted and applied to give the same performance guarantees.


\subsection{Computational Aspects}

Let us conclude by discussing computation aspects in the liquid democracy game.
Recall, to find an $\epsilon$-Nash equilibrium with a provably optimal social welfare guarantee
we solved a fixed point theorem. 
Moreover, solving the fixed point theorem requires knowledge of the optimal set $D^*$ of gurus.
But obtaining $D^*$ is equivalent to finding an optimal solution to the liquid democracy game 
and this problem is hard.
\begin{thm}\label{thm:hardness}
It is NP-hard to find an optimal solution to the liquid democracy game.
\end{thm}
\begin{proof}
We apply a reduction from {\em dominating set in a directed graph}. Given a directed 
graph $G=(V,A)$ and an integer $k$: is there a set $S\subseteq V$ of cardinality $k$ such that,
for each $i\notin S$, there exists an arc $(i,j)\in A$ for some $j\in S$. 
This problem is NP-complete. We now give a reduction to the liquid democracy game.

Given the directed graph $G=(V,A)$ let the set of agents in the game be $V$. Define the
utility function of an agent by:
$$
u_{ij} = 
\begin{cases}
1 &\text{if } (i,j)\in A  \\ 
0 & \text{otherwise}	
\end{cases} 
$$ 
It immediately follows that there is a dominating set of cardinality at most $k$ if and only if the
liquid democracy game has a solution with social welfare at least $n-k$.
This completes the proof.
\qed
\end{proof}

So is it possible for the agents to compute a good bicriteria solution in polynomial time? 
If we allow for sub-optimal approximation guarantees then this is achievable.
The idea is simple: the characteristics required of the agents to ensure 
reasonable performance guarantees are {\em narcissism} and {\em avarice}.\footnote{This is not an unreasonable assumption for
both pirates and many of the inhabitants of Wonderland!}
First, since $D^*$ is unknown, each agent $i$ narcissistically assumes he himself is an optimal guru.
Consequently, he will vote with probability $p$. Thus, he will delegate his vote 
with probability $(1-p)$. This he will do avariciously, by greedily delegating to the agent $i^*$
that gives him largest myopic utility, that is $i^*=\argmax_{j\in V} u_{ij}$.

\begin{thm}\label{thm:approx}
For any $\epsilon\in[\frac{3}{4},1]$, the narcissistic-avaricious algorithm is linear time and produces an
$\epsilon$-Nash equilibrium with social welfare at least $(1-\epsilon)\cdot \Opt$.
\end{thm}
\begin{proof}
Consider the incentive guarantee for agent $i$ for the strategy profile $\z$ induced by the algorithm.
As $i$ votes with probability $z_{ii}=p$ and delegates to $i^*$ with probability $z_{ii^*}=(1-p)$,
his utility is 
\begin{align*}
		u_i(\z_i,\z_{-i}) &\ \geq\ z_{ii}u_{ii} + z_{i i^*} z_{i^* i^*}  u_{ii^*} \\
		&\ =\ p\cdot u_{ii} + (1-p) \cdot p\cdot u_{ii^*} \\
		&\ \geq\ (1-p)\cdot p\cdot u_{ii^*} \\
		&\ \geq\ (1-p)\cdot p \cdot u_i(\hat{\z_i},\z_{-i}) \qquad \forall \hat{\z_i}\in \mathbb{S}_i
\end{align*}
Hence $\z$ is an $\epsilon$-Nash equilibrium for 
$$\epsilon=(1-p(1-p))=1-p+p^2=\frac34 +(p-\frac12)^2$$
It follows that the narcissistic-avaricious algorithm can provide incentive guarantees only for $\epsilon\in[\frac{3}{4},1]$.
Further, by solving the corresponding quadratic equation, to obtain such an $\epsilon$-Nash equilibrium we simply 
select $p=\frac12\left(1+\sqrt{1-4(1-\epsilon)}\right)$.

Now, let's evaluate the social welfare guarantee for the narcissistic-avaricious algorithm.	
As above, we have
 \begin{align*}
	 	\SW(\z) &\ =\ \sum_{i\in V} u_i(\z_i,\z_{-i})\\
		&\ \geq\ \sum_{i\in V} (1-p)\cdot p\cdot u_{ii^*} \\
	 	&\ =\  (1-p)\cdot p\cdot\sum_{i\in V} u_{ii^*}\\
		&\ \geq\ (1-p)\cdot p\cdot \Opt 
 \end{align*}
 Since $p=\frac12\left(1+\sqrt{1-4(1-\epsilon)}\right)$ this gives
 $$\SW(\z) \ \ge\ \frac14\Big(1-(1-4(1-\epsilon)\Big)\cdot \Opt \ = \ (1-\epsilon)\cdot \Opt$$
Therefore, as claimed, the narcissistic-avaricious algorithm outputs a solution whose welfare is at least $(1-\epsilon)$
times the optimal welfare. 

Finally, observe that implementing the narcissistic-avaricious strategy requires that each agent $i$ simply 
computes $i^*=\argmax_{j\in V} u_{ij}$. This can be done for every agent in linear time in the size of the input.
\qed
\end{proof}

We emphasize two points concerning Theorem~\ref{thm:approx}. One, it only works for weaker incentive guarantees, namely 
$\epsilon\in[\frac{3}{4},1]$. Unlike the fixed point algorithm it does not work for the range $\epsilon\in (0,\frac34)$.
Two, the social welfare guarantee is $(1-\epsilon)$. This is a constant but, for the valid range $\epsilon\in[\frac{3}{4},1]$, it is much worse 
than the $\epsilon$ guarantee obtained by the fixed point algorithm. 
A very interesting open problem is to find a polynomial time algorithm that matches the optimal bicriteria
guarantees provided by Theorem~\ref{thm:bicriteria} and applies for all $\epsilon>0$.

\vspace*{0.5cm}

\noindent \textbf{Acknowledgements}: We thank Bundit Laekhanukit for interesting discussions. We are also grateful to the reviewers for suggestions which greatly improved the paper.

